% chapters_parte2/06_kubernetes_networking.tex

\chapter{\eh{gear} Networking en Kubernetes}

Kubernetes es hoy la plataforma estándar de orquestación de contenedores. Su modelo
de red es intencionalmente simple en la interfaz pero complejo en la implementación:
\emph{todo Pod puede hablar con todo Pod sin NAT}. Entender cómo se implementa eso
es fundamental para operar clusters en producción.

% ─────────────────────────────────────────────
\section{\eh{memo} El Modelo de Red de Kubernetes}

\begin{concepto}
Kubernetes impone cuatro reglas fundamentales al networking:

\begin{enumerate}
    \item Cada Pod tiene su propia dirección IP única en el cluster.
    \item Todos los Pods pueden comunicarse entre sí sin NAT.
    \item Los nodos pueden comunicarse con todos los Pods sin NAT.
    \item La IP que ve un Pod de sí mismo es la misma que ven los demás Pods.
\end{enumerate}

Estas reglas simplifican el modelo: las aplicaciones no necesitan saber si están en el
mismo nodo o en nodos distintos. El networking subyacente lo resuelve.
\end{concepto}

\begin{cuidado}
Kubernetes \textbf{define} el modelo pero \textbf{no lo implementa}. La implementación
la provee el \textbf{CNI plugin} (Container Network Interface): Cilium, Calico, Flannel,
Weave, etc. Distintos CNIs tienen distintas capacidades (Network Policies, encryption,
eBPF, multicast, etc.).
\end{cuidado}

% ─────────────────────────────────────────────
\section{\eh{electric-plug} Tipos de red en un cluster}

\begin{center}
\begin{tikzpicture}[
    netbox/.style={draw=cyan!50!white, fill=darkcardgray, text=textlight,
                   rectangle, rounded corners=4pt, minimum width=5cm,
                   minimum height=0.9cm, font=\small\bfseries, align=center},
    node/.style={draw=green!50!white, fill=darkcardgreen, text=textlight,
                 rectangle, rounded corners=4pt, minimum width=3cm,
                 minimum height=0.8cm, font=\small, align=center},
    pod/.style={draw=gray!50!white, fill=darkcardgray, text=textlight,
                rectangle, minimum width=1.5cm, minimum height=0.6cm,
                font=\tiny, align=center},
    arrow/.style={->, thick, draw=gray!60!white},
]
    \node[netbox] (pod_net) at (0, 4.5) {Pod Network (ej. 10.244.0.0/16)};
    \node[netbox] (svc_net) at (0, 3)   {Service Network (ej. 10.96.0.0/12)};
    \node[netbox] (node_net) at (0, 1.5) {Node Network (ej. 192.168.1.0/24)};

    \node[node]  (n1) at (-3, 0) {Node 1\\192.168.1.10};
    \node[node]  (n2) at (3, 0)  {Node 2\\192.168.1.11};
    \node[pod]   (p1) at (-4.2, -1.5) {Pod A\\10.244.0.2};
    \node[pod]   (p2) at (-2.2, -1.5) {Pod B\\10.244.0.3};
    \node[pod]   (p3) at (1.8, -1.5)  {Pod C\\10.244.1.2};
    \node[pod]   (p4) at (3.8, -1.5)  {Pod D\\10.244.1.3};

    \draw[arrow] (n1) -- (p1); \draw[arrow] (n1) -- (p2);
    \draw[arrow] (n2) -- (p3); \draw[arrow] (n2) -- (p4);
    \draw[arrow, dashed, draw=yellow!60!white] (p1) to[bend left=30]
        node[font=\tiny, text=yellow!70!white, above] {CNI overlay} (p3);
\end{tikzpicture}
\end{center}

\begin{itemize}
    \item \textbf{Pod Network:} rango de IPs asignadas a Pods. CNI lo administra.
    \item \textbf{Service Network:} IPs virtuales de Services (ClusterIP). Solo existen
          en iptables/eBPF; ningún dispositivo real las tiene.
    \item \textbf{Node Network:} red física/real de los nodos. Preexiste a Kubernetes.
\end{itemize}

% ─────────────────────────────────────────────
\section{\eh{gear} Services y kube-proxy}

\begin{concepto}
Los Pods son efímeros: mueren y resurgen con nuevas IPs. Un \textbf{Service} es una
abstracción estable: tiene una IP virtual fija (ClusterIP) y un DNS name, y hace
load balancing hacia los Pods seleccionados por su \texttt{selector}.

\textbf{kube-proxy} implementa los Services programando reglas en el kernel del nodo:
\end{concepto}

\begin{center}
\begin{tabular}{lll}
\toprule
\textbf{Modo} & \textbf{Mecanismo} & \textbf{Rendimiento} \\
\midrule
userspace & Proxy en espacio de usuario & Lento, legado \\
iptables  & Reglas NAT en iptables       & Bueno, O(n) rules \\
IPVS      & IP Virtual Server del kernel & Mejor, O(1) lookup \\
eBPF      & Cilium sin kube-proxy        & Óptimo, bypass kernel stack \\
\bottomrule
\end{tabular}
\end{center}

\subsection{Tipos de Service}

\begin{itemize}
    \item \textbf{ClusterIP (default):} IP virtual solo accesible dentro del cluster.
          Uso: comunicación interna entre microservicios.
    \item \textbf{NodePort:} expone el Service en un puerto fijo de cada nodo
          (30000--32767). Accesible desde fuera del cluster via \texttt{NodeIP:NodePort}.
    \item \textbf{LoadBalancer:} solicita un load balancer externo al cloud provider
          (AWS ELB, GCP LB). Asigna IP pública automáticamente.
    \item \textbf{ExternalName:} alias DNS a un FQDN externo. Útil para integrar
          servicios externos como si fueran internos.
    \item \textbf{Headless (ClusterIP: None):} no hay IP virtual; DNS devuelve
          directamente las IPs de los Pods. Usado por StatefulSets (Kafka, Cassandra).
\end{itemize}

% ─────────────────────────────────────────────
\section{\eh{shield} Network Policies}

\begin{concepto}
Por defecto, todos los Pods en un cluster pueden comunicarse entre sí (modelo flat).
\textbf{NetworkPolicy} permite definir reglas de firewall a nivel de Pod: qué pods
pueden conectarse a qué otros pods, en qué puertos y protocolos.

Las NetworkPolicies son aditivas: si no hay política, se permite todo. Si hay al menos
una política que selecciona un Pod, ese Pod solo acepta lo que las políticas permiten.
\end{concepto}

\begin{codigo}
\begin{lstlisting}
# Solo permitir trafico al pod "backend" desde pods con label "tier: frontend"
apiVersion: networking.k8s.io/v1
kind: NetworkPolicy
metadata:
  name: backend-ingress
spec:
  podSelector:
    matchLabels:
      app: backend
  policyTypes:
  - Ingress
  ingress:
  - from:
    - podSelector:
        matchLabels:
          tier: frontend
    ports:
    - protocol: TCP
      port: 8080
\end{lstlisting}
\end{codigo}

\begin{cuidado}
NetworkPolicies requieren que el CNI plugin las soporte. Flannel no las implementa.
Calico, Cilium, y Weave sí. Sin un CNI compatible, los objetos NetworkPolicy se crean
pero no tienen efecto alguno --- la red sigue siendo completamente abierta.
\end{cuidado}

% ─────────────────────────────────────────────
\section{\eh{globe-with-meridians} Ingress y Gateway API}

\begin{concepto}
Los Services de tipo LoadBalancer son costosos (un LB externo por Service). \textbf{Ingress}
consolida el tráfico HTTP/HTTPS de múltiples Services detrás de un solo load balancer,
enrutando por host y path.
\end{concepto}

\begin{codigo}
\begin{lstlisting}
apiVersion: networking.k8s.io/v1
kind: Ingress
metadata:
  name: mi-app-ingress
spec:
  rules:
  - host: api.miempresa.com
    http:
      paths:
      - path: /v1
        pathType: Prefix
        backend:
          service:
            name: api-v1
            port:
              number: 80
      - path: /v2
        pathType: Prefix
        backend:
          service:
            name: api-v2
            port:
              number: 80
\end{lstlisting}
\end{codigo}

\begin{moderno}
\emoji{rocket} \textbf{Gateway API} (sucesor de Ingress): API más expresiva con soporte
nativo para TCP, UDP, gRPC, y modelos multi-tenant. Separa responsabilidades entre el
equipo de infraestructura (GatewayClass, Gateway) y el equipo de aplicación (HTTPRoute,
TCPRoute). Cilium, Istio y Nginx ya la implementan. Ingress quedará deprecated eventualmente.
\end{moderno}

% ─────────────────────────────────────────────
\section{\eh{satellite-antenna} CNI Plugins: Cilium vs Calico}

\begin{center}
\begin{tabular}{lll}
\toprule
\textbf{Feature} & \textbf{Cilium} & \textbf{Calico} \\
\midrule
Data plane         & eBPF                & iptables / eBPF (opcional) \\
Rendimiento        & Excelente           & Bueno \\
Network Policies   & L3/L4/L7 (HTTP, gRPC) & L3/L4 \\
Observabilidad     & Hubble (flujos)     & Limitada \\
Modo sin kube-proxy & Sí (eBPF)          & Parcial \\
BGP support        & Sí                  & Sí (más maduro) \\
Encryp. transparente & WireGuard/IPsec   & WireGuard/IPsec \\
\bottomrule
\end{tabular}
\end{center}

\begin{ejemplo}
\textbf{Por qué Cilium usa eBPF:}
Con iptables, cada Service agrega reglas que el kernel evalúa secuencialmente.
Con 10,000 Services, kube-proxy crea cientos de miles de reglas; la actualización
tarda minutos y la evaluación es O(n). eBPF usa mapas hash: lookup O(1), actualizaciones
atómicas en microsegundos. En clusters grandes ($>$500 nodos), Cilium reduce la latencia
de Service lookup de milisegundos a microsegundos.
\end{ejemplo}
